\documentclass[a4paper,12pt]{article}
\usepackage[utf8]{inputenc}
\usepackage[russian]{babel}
\usepackage{amsmath, amsfonts, amssymb}
\usepackage{geometry}
\usepackage{graphicx}
\usepackage{wrapfig}
\usepackage{tikz}

\geometry{left=25mm, right=20mm, top=20mm, bottom=20mm}

\begin{document}


\noindent 65] \hfill \S~3. КЛАССИФИКАЦИЯ БЕСКОНЕЧНО МАЛЫХ \hfill 163
\vspace{1em}
\hrule
\vspace{1.5em}

Для упрощения этой формулы вспомним, что
\begin{equation*}
    \alpha \doteq \sin \alpha = \frac{OD}{Oo} = \frac{R - r}{d}
\end{equation*}
--- в предположении, что $R - r$ мало относительно $d$. В том же предположении
\begin{equation*}
    \sqrt{d^2 - (R - r)^2} = d\sqrt{1 - \left(\frac{R - r}{d}\right)^2} \doteq d \left[ 1 - \frac{1}{2} \left( \frac{R - r}{d} \right)^2 \right].
\end{equation*}

После подстановки этих значений и преобразований, получим окончательную формулу:
\begin{equation*}
    l \doteq \pi(R + r) + 2d + \frac{(R - r)^2}{d}.
\end{equation*}

3) При разбивке дуг окружностей на местности имеет значение следующая задача: \textit{найти отношение стрелы $f = DB$ дуги $ABC$ окружности к стреле $f_1 = D_1B_1$ половины $AB_1B$ этой дуги (рис. 27).}

\begin{wrapfigure}{r}{0.4\textwidth}
    \centering
    \begin{tikzpicture}[scale=1.2]
    
        \coordinate (O) at (0,0);
        \coordinate (B) at (0,3);
        \coordinate (A) at (-2.2,2.04);
        \coordinate (C) at (2.2,2.04);
        \coordinate (D) at (0,2.04);
        
        \draw[thick] (A) arc (137:43:3);
        
        \draw (A) -- (C);
        \draw (O) -- (A);
        \draw (O) -- (C);
        \draw (O) -- (B);
        
        \coordinate (B1) at (-1.15, 2.77);
        \draw[dashed] (A) arc (137:112.5:3);
        \draw[dashed] (A) -- (B);
        
        \node[left] at (A) {$A$};
        \node[right] at (C) {$C$};
        \node[above] at (B) {$B$};
        \node[below right] at (D) {$D$};
        \node[below] at (O) {$O$};
        \node[above left] at (B1) {$B_1$};
        \node[below] at (-0.6, 2.4) {$D_1$};
        
        \draw[<->] (0.2, 2.04) -- (0.2, 3) node[midway, right] {$f$};
        \draw[<->, dashed] (-0.8, 2.5) -- (-0.4, 2.9) node[midway, above] {$f_1$};
        
        \node at (0.2, 0.8) {$\varphi$};
        \node at (-1.1, 1.5) {$r$};
    \end{tikzpicture}
    \small \textbf{Рис. 27}
\end{wrapfigure}

Если положить радиус окружности равным $r$, $\angle AOB = \varphi$, то $\angle AOB_1 = \frac{\varphi}{2}$ и
\begin{gather*}
    f = DB = r(1 - \cos \varphi), \\
    f_1 = r\left(1 - \cos \frac{\varphi}{2}\right).
\end{gather*}

Таким образом, искомое отношение равно
\begin{equation*}
    \frac{f}{f_1} = \frac{1 - \cos \varphi}{1 - \cos \frac{\varphi}{2}}.
\end{equation*}

Выражение это слишком сложно, чтобы им удобно было пользоваться на практике. Найдем его предел при $\varphi \to 0$ (ибо для достаточно малых $\varphi$ это выражение можно приближенно заменить его пределом). С этой целью заменяем числитель и знаменатель их главными частями и сразу находим:
\begin{equation*}
    \lim \frac{f}{f_1} = \lim \frac{\frac{1}{2}\varphi^2}{\frac{1}{2}\left(\frac{\varphi}{2}\right)^2} = 4.
\end{equation*}

Итак, для дуг, соответствующих небольшому центральному углу, приближенно можно считать, что \textit{стрела полудуги вчетверо меньше стрелы дуги}. Это позволяет последовательно строить промежуточные точки дуги, для которой даны концы и середина.

\vspace{1.5em}
\noindent \textbf{65. Классификация бесконечно больших.} Заметим, что для бесконечно больших величин может быть развита подобная же классификация. Как и в \textbf{60}, будем считать рассматриваемые бесконечно большие величины функциями от одной и той же переменной $x$, которые стремятся к $+\infty$, когда $x$ стремится к $a$.

\end{document}